\documentclass[10pt,a4paper]{amsart}
\usepackage[margin=19mm]{geometry}
\usepackage{amsmath,amssymb,mathtools}
\usepackage[scr=rsfs,cal=euler]{mathalfa}
\usepackage{xfrac}
\usepackage[unicode,hidelinks,bookmarks=false]{hyperref}
\newcommand{\ie}{i.e.\ }
\newcommand{\cf}{cf.\ }
\newcommand{\resp}{respectively\ }
\newcommand{\Subsection}{\subsection}
\providecommand{\nobreakdash}{\nobreak\hbox{-}}
\setlength{\emergencystretch}{2em}
\multlinegap0pt
\usepackage{bm}
\usepackage{bbm}
\usepackage{dsfont}
\usepackage{xspace}
\usepackage{cancel}
\usepackage{mathtools}
\usepackage[shortlabels]{enumitem}
\usepackage{amsthm}
\makeatletter
\@ifundefined{definition}{\newtheorem{definition}{Definition}}{}
\@ifundefined{theorem}{\newtheorem{theorem}{Theorem}}{}
\makeatother
\title[Generalized saddle-node ghosts and their composite structure]{Generalized saddle-node ghosts and their composite structures in dynamical systems}
\author{Daniel Koch, Akhilesh P. Nandan}
\date{Source: arXiv:2604.05194v3 (CC BY 4.0)}
\begin{document}
\maketitle

\noindent\emph{Source and licence.} Generalized saddle-node ghosts and their composite structures in dynamical systems. Version arXiv:2604.05194v3 (CC BY 4.0), distributed under the Creative Commons Attribution 4.0 International licence, as recorded in the arXiv licence metadata for this exact version and shown on the abstract page https://arxiv.org/abs/2604.05194v3. Full rights notice: licenses/arxiv-2604.05194-CC-BY-4.0.txt.\par\smallskip

\noindent\textit{Editorial scope.} Counted unit: the Theorem whose source label is \texttt{theorem\_jKSNinThetaNeurons} in the article (source0.tex lines 873--875, source label \texttt{theorem\_jKSNinThetaNeurons}), delivered together with the article's own complete original proof of it (source0.tex lines 877--966, one \texttt{proof} environment of 90 lines). No mathematics is written, repaired or re-derived: statement and proof are the article's own text line for line. Required context retained uncounted: def:SN (lines 168--176); eq:augustson (lines 278--281). Every definition, display and label of those lines is retained unchanged. Omitted from this package and not redistributed: 1--167, 177--277, 282--872, 876--876, 967--1130. Nothing inside a retained excerpt is omitted or altered: every retained body is delivered exactly as the article's lines are written, so no omission receipt of any kind accompanies this package. Citations: the retained excerpts carry no citation command at all, so no bibliography entry of the article is transcribed and no reference is redistributed. Reference adaptation: none was needed and none was made; every reference inside the delivered unit resolves inside it, so no unresolved reference or citation is produced. Printed slips kept literal: no printed word, symbol, hypothesis or number of the counted statement or of its proof was altered. Layout and class: the article's own class is \texttt{article} with packages including arxiv, amsmath, amsthm, inputenc, fontenc, hyperref, url, booktabs, amsfonts, nicefrac, microtype, lipsum, graphicx, enumitem; the wrapper is the lane's pinned \texttt{amsart} editorial wrapper at 10pt on A4, which restates the theorem-like environments the retained text needs and adds the article's own macro definitions that the retained text needs (none), with extra wrapper packages (bm, bbm, dsfont, xspace, cancel, mathtools, enumitem). The delivered numbering is the unit's own. Assets: no retained span contains a figure environment, a graphics-inclusion command, a PDF-page inclusion, a table or a TikZ picture; no image file is copied into this package, the package declares no asset dependency and no image file is redistributed with it. Licence: version arXiv:2604.05194v3 of this article is distributed under the Creative Commons Attribution 4.0 International licence, as recorded in the arXiv licence metadata for this exact version and shown on the abstract page https://arxiv.org/abs/2604.05194v3. A full rights notice is delivered at licenses/arxiv-2604.05194-CC-BY-4.0.txt. Characters: every retained excerpt is pure ASCII, so the delivered engine, XeLaTeX with its default Latin Modern text font, reports no missing character.
\begin{definition}[\textit{Saddle-node equilibrium and bifurcation point of type $j,k$}]\label{def:SN}
Let $\dot{x} = f(x,\rho_{\textnormal{crit}})$, $x\in \mathbb{R}^n$, $f\in \mathcal{C}^\omega$, be an autonomous dynamical system with parameters $\rho_{\textnormal{crit}} = (\rho_1,\dots,\rho_m)^T\in \mathbb{R}^m$. We say $x^*$ fulfilling $f(x^*)=0$ is a saddle-node equilibrium of type $j,k$ and $(x^*,\rho_{\textnormal{crit}})$ a saddle-node bifurcation point of type $j,k$ if the following holds true:
\begin{enumerate}[label={(\textbf{SN\arabic*})}, ref={(\textbf{G\arabic*})}]
    \item $D_xf(x^*,\rho_{\textnormal{crit}})$ has a semi-simple eigenvalue 0 of algebraic and geometric multiplicity $j$, with corresponding linearly independent eigenvectors  $v_1,\dots,v_j$ and $w_1,\dots,w_j$ to the right and left, respectively, that vary continuously with $(\rho_1,\dots,\rho_m)^T$.
    \item For each $i$ with $1\leq i \leq j$ there exists $l$ with $1\leq l \leq m$ with $w_i((\partial f/\partial\rho_l)(x^*,\rho_{\textnormal{crit}})) \neq 0$.
    \item for every non-zero $u\in \ker(D_xf(x^*,\rho_{\textnormal{crit}}))$, there exists $i$ with $1\leq i \leq j$ such that \mbox{$w_i(D_x^2f(x^*,\rho_{\textnormal{crit}})(u,u))\neq 0$}.
    \item $D_xf(x^*,\rho_{\textnormal{crit}})$ has $k$ eigenvalues $\lambda^+_1,\dots,\lambda^+_k$ with positive real parts, and $n-j-k$ eigenvalues $\lambda^-_1,\dots,\lambda^-_{n-j-k}$ with negative real parts.
\end{enumerate}
\end{definition}

\begin{align}
    \dot{\theta}_1 &= 1 - \cos(\theta_1) + \left(1 + \cos(\theta_1)\right) \left(\eta_1 + \kappa \left(1 - \cos(\theta_2)\right)\right) \label{eq:augustson}\\
    \dot{\theta}_2 &= 1 - \cos(\theta_2) + \left(1 + \cos(\theta_2)\right) \left(\eta_2 + \kappa \left(1 - \cos(\theta_1)\right)\right) \nonumber,
\end{align}

\begin{theorem}[Existence of a type $2,0$ saddle-node equilibrium in coupled theta neurons] \label{theorem_jKSNinThetaNeurons}
For $\eta_1 = \eta_2  = 0$ and $\kappa >0$, system \ref{eq:augustson} has a type $2,0$ saddle-node equilibrium at the origin.
\end{theorem}

\begin{proof} 
Denote $\mathbf{\theta} = (\theta_1,\theta_2)$. For $\theta=(0,0)$ we have 

\begin{align*}
\dot{\theta}_1 &= 1 - \cos\theta_1 + \left(1 + \cos\theta_1\right) \left(\eta_1 + \kappa \left(1 - \cos\theta_2\right)\right)\\ 
&= 1 - \cos0 + \left(1 + \cos0\right) \left(0 + \kappa \left(1 - \cos0\right)\right)\\
&= 1 - 1 + \left(1 + 1\right) \left(0 + \kappa \left(1 - 1\right)\right) = 0\\
\end{align*}
and analogously $\dot \theta_2=0$, showing that the origin is an equilibrium point.
We next prove that the origin fulfills (SN1) from definition \ref{def:SN}. The Jacobian matrix of system \ref{eq:augustson} is
$$D_{\theta} f(\theta,(\eta_1,\eta_2,\kappa))=\begin{bmatrix}
\sin\theta_1(1-\eta_1-\kappa(1-\cos\theta_2)) & \kappa\sin\theta_2(1+\cos\theta_1) \\
\kappa\sin\theta_1(1+\cos\theta_2) & \sin\theta_2(1-\eta_2-\kappa(1-\cos\theta_1)) 
\end{bmatrix}.$$ We thus have  
$$D_{\theta} f((0,0),(0,0,\kappa))=\begin{bmatrix}
0 & 0 \\
0 & 0
\end{bmatrix},$$ making the equilibrium at the origin non-hyperbolic with eigenvalues $\lambda_1=\lambda_2=0$. That is, have a geometric multiplicity of $2$ for the eigenvalue $0$. To find the eigenvectors $w\in \mathbb{R}^2$, we need to solve $$(D_{\theta} f((0,0),(0,0,\kappa))-\lambda_i I)v=
\begin{bmatrix}
-\lambda_i & 0 \\
0 & -\lambda_i
\end{bmatrix}
\begin{pmatrix}
v_1 \\
v_2
\end{pmatrix}
= \begin{pmatrix} 0 \\ 0\end{pmatrix}.
$$
Since $\lambda_i=0$, every non-zero vector $v\in\mathbb{R}^2$ fulfills above equation and is thus an eigenvector of $\lambda_i$. Therefore, the center eigenspace is $E^c=\{v: (D_{\theta} f((0,0),(0,0,\kappa))-\lambda_i I) \ v=0\} = \mathbb{R}^2$. Since $\textnormal{dim}(E^c)=2$, we have a geometric multiplicity of $2$, making $\lambda_i$ semi-simple with eigenvalue $v_i$ to the right. Since the Jacobian is symmetric, we can simply set $w_i=v_i^T$ as the eigenvector to the left. Hence (SN1) is satisfied.\\ \\
We next show that (SN2) is fulfilled. Denote $$f((\theta_1,\theta_2),(\eta_1,\eta_2,\kappa)) = \begin{pmatrix}  1 - \cos\theta_1 + \left(1 + \cos\theta_1\right) \left(\eta_1 + \kappa \left(1 - \cos\theta_2\right)\right) \\  1 - \cos\theta_2 + \left(1 + \cos\theta_2\right) \left(\eta_2 + \kappa \left(1 - \cos\theta_1\right)\right)\end{pmatrix}.$$
Then $\frac{\partial f}{\partial \eta_1}((\theta_1,\theta_2),(\eta_1,\eta_2,\kappa)) = \begin{pmatrix}1 + \cos\theta_1 \\ 0\end{pmatrix}$, so $w_i\frac{\partial f}{\partial \eta_1}((0,0),(0,0,\kappa))=w_i\begin{pmatrix}2\\0\end{pmatrix}$. Analogously we have $\frac{\partial f}{\partial \eta_2}((\theta_1,\theta_2),(\eta_1,\eta_2,\kappa)) = \begin{pmatrix}0\\1 + \cos\theta_2\end{pmatrix}$, so $w_i\frac{\partial f}{\partial \eta_2}((0,0),(0,0,\kappa))=w_i\begin{pmatrix}0\\2\end{pmatrix}$. Since $w_i\neq\begin{pmatrix}0 \ 0\end{pmatrix}$, we thus have $w_i\frac{\partial f}{\partial \eta_n}((0,0),(0,0,\kappa))\neq0$ for either $n=1$ or $n=2$, showing that (SN2) is fulfilled.
\\ \\
We now show that (SN3) is fulfilled, too. Denote $f_1=\dot\theta_1$ and $f_2=\dot\theta_2$, giving the the following second order partial derivatives:

\begin{align*}
    \frac{\partial^2 f_1}{\partial^2 \theta_1} &= \cos\theta_1(1-\eta_1-\kappa(1-\cos\theta_2)),\quad
    \frac{\partial^2 f_1}{\partial \theta_1 \partial \theta_2} = -\sin\theta_1\sin\theta_2,\quad
    \frac{\partial^2 f_1}{\partial^2 \theta_2} = \kappa\cos\theta_2(1-\cos\theta_1),\\
    \frac{\partial^2 f_2}{\partial^2 \theta_1} &= \kappa\cos\theta_1(1-\cos\theta_2),\quad
    \frac{\partial^2 f_2}{\partial \theta_1 \partial \theta_2} = -\sin\theta_1\sin\theta_2, \quad
    \frac{\partial^2 f_"}{\partial^2 \theta_2} =  \cos\theta_2(1-\eta_2-\kappa(1-\cos\theta_1)).
\end{align*}
 The Hessians are thus:

 $$
 H_1 = \begin{bmatrix}
\frac{\partial^2 f_1}{\partial^2 \theta_1} & \frac{\partial^2 f_1}{\partial \theta_1 \partial \theta_2}
\\ \frac{\partial^2 f_1}{\partial \theta_1 \partial \theta_2} & \frac{\partial^2 f_1}{\partial^2 \theta_2} 
 \end{bmatrix}
 =
\begin{bmatrix}
\cos\theta_1(1-\eta_1-\kappa(1-\cos\theta_2)) &  -\sin\theta_1\sin\theta_2 \\
 -\sin\theta_1\sin\theta_2 &\kappa\cos\theta_2(1-\cos\theta_1) 
 \end{bmatrix},$$ so $H_1|_{(\theta_1,\theta_2)=(0,0),\eta_1=\eta_2=0}=
 \begin{bmatrix}
     1 & 0 \\ 0 & 2\kappa \\
 \end{bmatrix}.$
 Analogously, we get to
 $$
 H_2 = \begin{bmatrix}
\frac{\partial^2 f_1}{\partial^2 \theta_1} & \frac{\partial^2 f_2}{\partial \theta_1 \partial \theta_2}
\\ \frac{\partial^2 f_2}{\partial \theta_1 \partial \theta_2} & \frac{\partial^2 f_2}{\partial^2 \theta_2} 
 \end{bmatrix}
 =
\begin{bmatrix}
\kappa\cos\theta_1(1-\cos\theta_2) &  -\sin\theta_1\sin\theta_2 \\
 -\sin\theta_1\sin\theta_2 &\cos\theta_2(1-\eta_2-\kappa(1-\cos\theta_1)) 
 \end{bmatrix},$$ so $H_2|_{(\theta_1,\theta_2)=(0,0),\eta_1=\eta_2=0} =
 \begin{bmatrix}
     2\kappa & 0 \\ 0 & 1 \\
 \end{bmatrix}.$ 
\\ Now let $u = \begin{pmatrix}a\\b \end{pmatrix}\in \ker(D_{\theta} f((0,0),(0,0,\kappa)))$, then $u^TH_1u=(a\ b)\left(\begin{pmatrix}
     1 & 0 \\ 0 & 2\kappa \\
 \end{pmatrix}\begin{pmatrix}
     a\\b
 \end{pmatrix}\right)=a^2+2\kappa b^2$ and
 $u^TH_2u=(a\ b)\left(\begin{bmatrix}
     2\kappa & 0 \\ 0 & 1 \\
 \end{bmatrix}\begin{pmatrix}
     a\\b
 \end{pmatrix}\right)=2\kappa a^2+b^2$. Hence $D_\theta^2f((0,0),(0,0,\kappa))(u,u)=\begin{pmatrix}
     a^2+2\kappa b^2\\2\kappa a^2+b^2
 \end{pmatrix}.$
 

Suppose $w_iD_\theta^2f((0,0),(0,0,\kappa))(v,v)= c(a^2+2\kappa b^2)+d(2\kappa a^2+b^2)=0$
for all non-zero $w_i = (c \ d)$, which is only possible if $a^2+2\kappa b^2=0$ and $2\kappa a^2+b^2=0$. Since $\kappa>0$, it follows that $a=b=0$, so $v=\begin{pmatrix} 0\\0 \end{pmatrix}$. Thus if $v = \begin{pmatrix}a\\b \end{pmatrix}\in \ker(D_{\theta} f((0,0),(0,0,\kappa)))$, by contraposition it follows that if $v \neq \begin{pmatrix} 0\\0\end{pmatrix}$, then $w_iD_\theta^2f((0,0),(0,0,\kappa))(v,v)\neq 0$ for at least one $1\leq i\leq j$, satisfying (SN3).\\

 Finally, since we showed that $D_{\theta} f((0,0),(0,0,\kappa))$ has no other eigenvalues than $0$, it is $k=0$ (SN4), completing the proof.
\end{proof}

\end{document}
